% ECONOMICS_PROBLEM: {"id": "E62-167", "title": "Effective fiscal-rule design", "jel_code": "E62", "source_id": "OP-0259"}
% !TeX program = xelatex
\documentclass[11pt,a4paper]{article}
\usepackage[margin=23mm]{geometry}
\usepackage{fontspec}
\setmainfont{Latin Modern Roman}
\usepackage{amsmath,amssymb,mathtools}
\usepackage{xcolor,enumitem,booktabs,longtable,array,xurl,hyperref}
\definecolor{muted}{HTML}{53636C}
\hypersetup{colorlinks=true,urlcolor=blue,linkcolor=blue}
\setlength{\parindent}{0pt}
\setlength{\parskip}{5pt}
\newcommand{\R}{\mathbb R}
\newcommand{\E}{\mathbb E}
\newcommand{\N}{\mathbb N}
\newcommand{\ind}{\mathbf 1}
\newcommand{\Var}{\operatorname{Var}}
\newcommand{\Cov}{\operatorname{Cov}}
\newcommand{\supp}{\operatorname{supp}}
\DeclareMathOperator*{\argmin}{arg\,min}
\DeclareMathOperator*{\argmax}{arg\,max}
\newcommand{\entrymeta}[3]{{\sffamily\small\color{muted}\textbf{\texttt{#1}}\quad|\quad #2\par #3\par}}
\newcommand{\Problemref}[1]{\texttt{#1}}
\newcommand{\JELref}[2]{\texttt{#2} (JEL \texttt{#1})}
\begin{document}
\section*{Effective fiscal-rule design}
\textbf{Problem ID:} \texttt{E62-167}\quad \textbf{JEL:} \texttt{E62}\par
% BEGIN ECONOMICS STATEMENT
\label{problem:OP-0259}
{\sffamily\small\color{muted}Original subject group: Fiscal policy and sovereign debt\par}
\label{definitions:problem:30007540}
\textbf{Audit classification:} Background; proposed extension. The source recommends fiscal-rule improvements; a globally optimal incentive-compatible rule is not established as open. Verification concerns the cited version, not first priority or proof that this exact editorial specification remains unsolved on October 8, 2026.
\subsection*{Definitions and precise research target}
\paragraph{Editorial mathematical specification.} Fix a stochastic fiscal economy with debt ratio \(b_t\), output growth \(g_t\), real financing rate \(r_t\), and primary surplus \(s_t\), expressed as a ratio to output at \(t+1\), with \(g_t>-1\). Debt evolves as \(b_{t+1}=(1+r_t)b_t/(1+g_t)-s_t\). A political government values resident welfare plus a declared electoral benefit from current spending. Fiscal institutions select rule parameters \(a=(\phi,\bar b,e,p)\): a debt-feedback coefficient, reference debt ratio, measurable escape-clause trigger, and enforcement probability or penalty schedule. The government retains explicitly bounded discretion when the trigger activates.
For a stated admissible rule family \(\mathcal A\), identify the rule maximizing expected resident welfare subject to budget feasibility, government incentive compatibility, and a bound on crisis risk. Formally,
\[
a^*\in\arg\max_{a\in\mathcal A}E_a\sum_{t=0}^{H}\beta^t u(c_t,\ell_t).
\]
Allow compliance, attempted circumvention, and rule revisions to be endogenous rather than assuming the announced rule is always followed. Estimate the response of those behaviors to monitoring and enforcement using credible institutional variation, and quantify the insurance value of escape clauses against their opportunistic use. Compare alternative rule designs under common shocks and political preferences. This is a proposed institutional-design problem; successful empirical associations between rules and debt performance do not establish an optimal universally applicable rule.

Require \(r_t>-1\) and define \(g_t=Y_{t+1}/Y_t-1\) for positive real GDP; \(s_t=S_{t+1}/Y_{t+1}\) is the real primary surplus in the next-period output denominator. Give a rule such as \(s_t\ge\phi(b_t-\bar b)\) outside a measurable escape event, with bounded relaxation on that event, and a specified monitoring/penalty technology. A political government chooses feasible policy to maximize resident utility plus its declared electoral payoff; the designer anticipates that response. Government incentive compatibility means prescribed compliance is at least as valuable as each permitted deviation, including the probability and cost of enforcement. If several government responses are equilibria, fix a selection or optimize worst-case resident welfare. This institutional-control formulation contains strategic incentives and is not wholly outside game theory.
\subsection*{Variants and assumption changes}
The following are \emph{editorial variants}, not additional published open conjectures.
\begin{enumerate}
\item \emph{Exogenous compliance.} Fix implementation probability and policy reactions outside escape episodes. Optimize the debt-feedback coefficient subject to fiscal feasibility and crisis probability; this benchmark deliberately abstracts from strategic evasion.
\item \emph{Endogenous evasion.} Add observable monitoring and costly concealment of off-budget liabilities. Choose enforcement and escape-clause design jointly under an administrative budget and compare with a rule that measures only headline debt.
\end{enumerate}
\subsection*{References and source audit}
\begin{itemize}
\item Julien Acalin; Virginia Alonso-Albarran; Clara Arroyo; Waikei Raphael Lam; Leonardo Martinez; Anh D. M. Nguyen; Francisco Roch; Galen Sher; Alexandra Solovyeva. \emph{Fiscal Guardrails against High Debt and Looming Spending Pressures}. 2025. Locator: Official SDN2025/004 record, Main findings and Policy implications; publication date September 25, 2025. Primary record and abstract/summary checked; no fulltext claim is made. \url{https://www.imf.org/en/publications/staff-discussion-notes/issues/2025/09/22/fiscal-guardrails-against-high-debt-and-looming-spending-pressures-569841}.
\end{itemize}
The source recommends fiscal-rule improvements; a globally optimal incentive-compatible rule is not established as open. This formulation contains strategic incentives and therefore overlaps game theory.
\begin{enumerate}\raggedright
\item\hypertarget{definitions:source:30007540:1}{} Julien Acalin; Virginia Alonso-Albarran; Clara Arroyo; Waikei Raphael Lam; Leonardo Martinez; Anh D. M. Nguyen; Francisco Roch; Galen Sher; Alexandra Solovyeva. 2025. \emph{Fiscal Guardrails against High Debt and Looming Spending Pressures}. Official SDN2025/004 record, Main findings and Policy implications; publication date September25,2025.
\url{https://www.imf.org/en/publications/staff-discussion-notes/issues/2025/09/22/fiscal-guardrails-against-high-debt-and-looming-spending-pressures-569841}.
DOI: \url{https://doi.org/10.5089/9798229023801.006}.
\end{enumerate}

\subsection*{Common conventions: Source mathematical conventions}

These conventions apply to each entry unless explicitly replaced.

\textbf{Models and identification.} A primitive model \(m\) specifies all probability laws, feasible choices, information, equilibrium and measurement rules used by its target. The admissible class \(\mathcal M\) is fixed independently of outcomes. If \(P_O(m)\) is its observable probability law and \(\tau(m)\) the target, its identified set at observable law \(P\) is
\[
 \mathcal I_\tau(P)=\{\tau(m):m\in\mathcal M,\ P_O(m)=P\}.
\]
Point identification means that this set is a singleton. An empty set signals incompatibility of data and assumptions. Coordinate bounds need not describe the joint set. Bounds are sharp only if every retained target is attainable or arbitrarily approachable by compatible models. Finite bounds require explicit restrictions on unobserved quantities; unrestricted confounding, hidden wealth, extrapolation or arbitrary equilibrium selection can yield uninformative sets.

\textbf{Causal interventions.} A potential outcome \(Y_i(p)\) is the outcome under a complete policy regime \(p\), including timing, financing, information and market spillovers. The target population and units are fixed before treatment. With interference, use the complete assignment vector or a justified exposure mapping; individual no-interference notation alone cannot identify an economy-wide intervention. A difference of mean potential outcomes does not require the cross-world joint distribution. Conditioning on post-treatment migration, survival, enrollment or employment changes the estimand and needs separate justification.

\textbf{Statistical statements.} An estimator, confidence set, forecast or test specifies a sampling law, model class and loss/coverage criterion. Uniform \((1-\alpha)\)-coverage means \(\inf_{m\in\mathcal M}\Pr_m\{\tau(m)\in C_n\}\geq1-\alpha\), or an explicitly asymptotic version. The whole parameter space trivially has coverage, so informativeness requires a stated width, risk or power objective. A forecasting improvement is relative to a named benchmark, information set and evaluation population.

\textbf{Equilibrium and optimization.} An equilibrium comprises feasible plans, optimality, beliefs where needed, budget constraints and all market/resource clearing. The primitives or a stated benchmark define each correspondence \(\mathcal E(p;m)\). Nonemptiness is assumed or proved, never inferred just from compact policy sets. A selected-equilibrium target uses a declared selection rule; a robust target quantifies over all permitted equilibria. Use suprema or infima unless attainment is established. An \(\varepsilon\)-optimal policy is feasible and within \(\varepsilon\) of the finite optimum value. Report an empty feasible set separately.

\textbf{Welfare and accounting.} Normative utility, interpersonal normalization, social weights, numeraire, discounting and the use of public receipts are primitives. Behavioral utility need not coincide with the welfare criterion. Household welfare includes ownership income and rents once; adding profits again to compensating variation generally double-counts them. A monetary equivalent is defined by an explicit transfer and price convention, with monotonicity and range conditions. Average effects, mechanism contributions, transfers and welfare losses are different targets. Infinite sums require convergence and dynamic feasible plans require terminal or no-Ponzi conditions.

\textbf{Source versions.} A working-paper date, revision date and journal publication year are distinguished. A DOI for a journal version is not automatically the DOI of an earlier manuscript. Locators refer to the stated version and printed pages where specified. A metadata-only check does not verify a claim about a full-text conclusion. Every entry's source subsection narrows the provenance claim accordingly.
% END ECONOMICS STATEMENT
\end{document}
